\section{The fixed theta denominator and the passage to short generators}
\label{sec:ftb-denominator-generators}

This section follows the frozen source at
\srcid{7cbf7a0c32fca1badf05931879a37c3f8a2b03b4}.
The numerical estimates and the lattice argument below are independent
of the analytic theta spanning theorem.  The latter is used only
through the explicit spanning and pointwise norm statements recorded
at the end of this section.

Let \(E\) be the CM host, \(\Psi\) its chosen CM type,
\(f=[E:\mathbb Q]/2\), and let the signature be \((p,q)\), with
\(r=p+q\).  Let \(K\) be the coefficient number field, \(M\ge3\) the odd
level, \(n\) an even positive integer, and
\(a=(a_{\tau,i})\) the nonnegative integral row degrees of a harmonic
polynomial.  Put
\[
 d_a=\sum_{\tau:E\hookrightarrow\mathbb C}
             \sum_{0\le i<r_\tau}a_{\tau,i},\qquad
 \kappa=d_a+fr\,\frac n2,\qquad
 \Delta_E=|\operatorname{disc}E|.
\]
Here \(r_\tau=p\) on \(\Psi\) and \(r_\tau=q\) on its conjugate.
The corresponding algebraic weight is
\[
 \lambda_{\tau,k}=
 \begin{cases}
 a_{\tau,k}+n/2,&0\le k<p,\\
 -a_{\bar\tau,r-1-k}-n/2,&p\le k<r .
 \end{cases}
\]
Thus \(\kappa=\sum_{\tau\in\Psi,k}|\lambda_{\tau,k}|\).
The source writes these quantities as
\srcid{DF}, \srcid{sizeK}, and \srcid{laOf}.

The exact integer used for all theta vectors with these row degrees is
\begin{equation}\label{eq:ftb-intZ}
 Z=M^{frn+d_a}\Delta_E^{d_a},\qquad N_a=Z^2.
\end{equation}
The factor \(M^{frn}\) is the cardinality of the finite Fourier-label
group \((\operatorname{LevQ})^n\).  In particular \(Z\) is chosen before
the Fourier label, the auxiliary Hermitian lattice, its denominator,
the harmonic polynomial, the vector represented by the theta rule,
or the determinant-balancing scalar.

\begin{srlemma}[The numerical denominator bound]\label{lem:ftb-intZ-size}
Suppose \(f,r,L\ge1\),
\(\log M\le30L\), and \(\log\Delta_E\le7fL\).
Then \(Z>0\) and
\[
 \log Z\le67(frL)^2\kappa,\qquad
 \log N_a\le134(frL)^2\kappa .
\]
In the source's good-cover packet the discriminant inequality follows
from \(\operatorname{fieldSz}(E)\le3L\) and \([E:\mathbb Q]=2f\);
the bound \(7fL\) is retained exactly as used there.
\end{srlemma}
\begin{proof}
The discriminant and \(M\) are nonzero integers, so \(Z>0\).
Evenness of \(n\) and nonnegativity of the row degrees give
\(d_a\le\kappa\) and \(frn+d_a\le2\kappa\).  Consequently
\[
 \log Z=(frn+d_a)\log M+d_a\log\Delta_E
 \le60L\kappa+7fL\kappa
 \le67fL\kappa
 \le67(frL)^2\kappa .
\]
The last inequality uses \(f,r,L\ge1\).  Since \(N_a=Z^2\), its logarithm
is twice that of \(Z\).  This is the argument in
\srcid{LogShortThetaSeriesDenominatorBoundLe.intZ\_log\_le}.
\end{proof}

We next give in full the arithmetic lattice step used after the theta
spanning family has been made integral.

\begin{srlemma}[The source's Euclidean transference estimate]
\label{lem:ftb-euclidean-transference}
Let \(\Lambda\) be a full lattice in a Euclidean space of dimension
\(d\).  Suppose every nonzero member of \(\Lambda\) has length at least
\(\lambda>0\).  Its Euclidean dual contains \(d\) linearly independent
vectors, each of length at most \(d^2/\lambda\).
\end{srlemma}
\begin{proof}
First the dual has a nonzero vector of length at most \(d/\lambda\).
In orthonormal coordinates, Minkowski's theorem for cubes says that a
lattice with covolume \(D<t^d\) has a nonzero vector of coordinate
supremum norm less than \(t\); the weak inequality follows by applying
the closed-cube version.  Since Euclidean length is at most
\(\sqrt d\) times the coordinate supremum norm, the minimum hypothesis
implies
\[
 \operatorname{covol}(\Lambda)\ge(\lambda/\sqrt d)^d.
\]
The dual has reciprocal covolume, so the closed-cube version supplies a
nonzero dual vector with coordinate norm at most \(\sqrt d/\lambda\),
hence Euclidean length at most \(d/\lambda\).

We prove the asserted independent-family estimate by induction on
\(d\); the zero-dimensional assertion is empty.  Divide the nonzero
dual vector just obtained by the positive greatest common divisor of
its integral pairings with a lattice basis.  The result \(y_0\) remains
in the dual, has no larger length, and has a lattice vector \(x_u\)
with \(\langle y_0,x_u\rangle=1\).
The homomorphism \(v\mapsto\langle y_0,v\rangle\) maps \(\Lambda\) onto
\(\mathbb Z\).  Its kernel
\(\Lambda_W=\Lambda\cap y_0^\perp\) is a full lattice in \(y_0^\perp\):
indeed \(v-\langle y_0,v\rangle x_u\) lies in this kernel, and these
vectors, as \(v\) runs through a basis of \(\Lambda\), span
\(y_0^\perp\).  Its nonzero lengths are still at least \(\lambda\).

Write \(d=m+1\).  The induction hypothesis gives independent vectors
\(w_1,\ldots,w_m\) in the dual of \(\Lambda_W\), of length at most
\(m^2/\lambda\).  Regard them as vectors in \(y_0^\perp\), and put
\[
 u_i=w_i-
  \bigl(\langle w_i,x_u\rangle-
             \operatorname{round}\langle w_i,x_u\rangle\bigr)y_0 .
\]
For \(v\in\Lambda\), let \(k=\langle y_0,v\rangle\in\mathbb Z\) and
\(v_0=v-kx_u\in\Lambda_W\).  Then
\[
 \langle u_i,v\rangle
 =\langle w_i,v_0\rangle+
           k\,\operatorname{round}\langle w_i,x_u\rangle\in\mathbb Z.
\]
Thus \(u_i\) belongs to the full dual.  Orthogonal projection onto
\(y_0^\perp\) proves that \(y_0,u_1,\ldots,u_m\) are independent.
Since the rounding error has absolute value at most \(1/2\),
\[
 \|u_i\|^2
 \le \frac{m^4+(m+1)^2/4}{\lambda^2}
 \le \frac{(m+1)^4}{\lambda^2}.
\]
Also \(\|y_0\|\le(m+1)/\lambda\le(m+1)^2/\lambda\).
This proves the induction.  These are precisely the cube,
primitivity, hyperplane-lattice, and rounding arguments in
\srcid{TransferenceA}, \srcid{TransferenceB}, and
\srcid{Transference.gtE\_main}.
\end{proof}

\begin{srlemma}[Short independent algebraic integers]
\label{lem:ftb-short-integers}
Let \(K\) be a number field of degree \(d\), with
\(\Delta_K=|\operatorname{disc}K|\).
There are algebraic integers \(\omega_1,\ldots,\omega_d\), linearly
independent over \(\mathbb Q\), such that
\[
 |\sigma(\omega_i)|\le d^3\Delta_K^{1/d}
 \quad\hbox{for every }\sigma:K\hookrightarrow\mathbb C
 \hbox{ and every }i .
\]
\end{srlemma}
\begin{proof}
List all \(d\) complex embeddings, including both members of every
nonreal conjugate pair, and put
\[
 \Phi_\pm(x)=
       \bigl(\operatorname{Re}\sigma(x)
                   \pm\operatorname{Im}\sigma(x)\bigr)_\sigma .
\]
Conjugate terms cancel the mixed products, so
\[
 \langle\Phi_+(x),\Phi_-(y)\rangle
       =\operatorname{Tr}_{K/\mathbb Q}(xy),\qquad
 \|\Phi_\pm(x)\|^2=\sum_\sigma|\sigma(x)|^2 .
\]
The trace form is nondegenerate; these maps identify the real
extensions of \(K\) with \(\mathbb R^d\).
An integral basis and its trace-dual basis therefore identify the
Euclidean dual of
\(\Phi_-(\mathfrak D_K^{-1})\) with \(\Phi_+(\mathcal O_K)\).

For a nonzero \(x\in\mathfrak D_K^{-1}\), the integral ideal
\((x)\mathfrak D_K\) has norm at least one.  Hence
\(|N_{K/\mathbb Q}x|\ge\Delta_K^{-1}\).
Applying arithmetic--geometric mean to
\((|\sigma(x)|^2)_\sigma\) gives
\[
 \sum_\sigma|\sigma(x)|^2
       \ge d\,\Delta_K^{-2/d}\ge\Delta_K^{-2/d}.
\]
Thus the nonzero vectors of
\(\Phi_-(\mathfrak D_K^{-1})\) have length at least
\(\Delta_K^{-1/d}\).  Lemma~\ref{lem:ftb-euclidean-transference}
supplies \(d\) independent vectors in its dual with lengths at most
\(d^2\Delta_K^{1/d}\).  They are \(\Phi_+(\omega_i)\) for
\(\omega_i\in\mathcal O_K\).
Each \(|\sigma(\omega_i)|\) is at most
\(\|\Phi_+(\omega_i)\|\), and \(d^2\le d^3\).
This proves the source's stated, deliberately weaker, bound.
\end{proof}

\begin{srproposition}[A bounded integral spanning set gives short
\(\mathbb Z\)-generators]\label{prop:ftb-short-generators}
Let \(V\) be a finite-dimensional \(K\)-vector space, and \(L\subset V\)
a finitely generated \(\mathcal O_K\)-submodule.
For every embedding \(\sigma:K\hookrightarrow\mathbb C\), let
\(N_\sigma:V\to\mathbb R_{\ge0}\) satisfy
\[
 N_\sigma(av)=|\sigma(a)|N_\sigma(v),\qquad
 N_\sigma(v+w)\le N_\sigma(v)+N_\sigma(w).
\]
Suppose a subset \(S\subset L\) spans \(V\) over \(K\), and
\(N_\sigma(s)\le A\) for every \(s\in S\) and every \(\sigma\).
Put \(m=\dim_KV\), \(d=[K:\mathbb Q]\).
Then \(L\) has at most \(2dm\) generators as a \(\mathbb Z\)-module,
all of whose \(N_\sigma\)-norms are at most
\[
 A_{\rm gen}=\max\{1,d^4m\Delta_K^{1/d}A\}.
\]
In particular \(A_{\rm gen}\ge1\).
\end{srproposition}
\begin{proof}
Choose a \(K\)-basis \(v_1,\ldots,v_m\) from \(S\), and the independent
integers of Lemma~\ref{lem:ftb-short-integers}.
Then \(E_{ij}=\omega_i v_j\) is a \(\mathbb Q\)-basis of \(V\), lies in
\(L\), and has all its norms at most
\(C=d^3\Delta_K^{1/d}A\).

For \(x\in V\), write \(x=\sum_{i,j}q_{ij}E_{ij}\) with
\(q_{ij}\in\mathbb Q\), and define
\[
 \operatorname{red}(x)
       =\sum_{i,j}(q_{ij}-\lfloor q_{ij}\rfloor)E_{ij}.
\]
Then
\(x-\operatorname{red}(x)\in\sum_{i,j}\mathbb Z E_{ij}\), while the
triangle inequality and \(0\le q-\lfloor q\rfloor<1\) give
\[
 N_\sigma(\operatorname{red}(x))\le dmC .
\]
As a finite torsion-free \(\mathbb Z\)-module, \(L\) has a
\(\mathbb Z\)-basis \(b_1,\ldots,b_t\), with \(t\le dm\):
its members are independent over \(\mathbb Q\) in \(V\).
Every \(\operatorname{red}(b_\nu)\) belongs to \(L\), since \(b_\nu\)
and \(b_\nu-\operatorname{red}(b_\nu)\) do.
The \(dm+t\) vectors
\[
 \{E_{ij}\}_{i,j}\ \cup\
       \{\operatorname{red}(b_\nu)\}_{\nu=1}^t
\]
generate \(L\) over \(\mathbb Z\), because they generate each \(b_\nu\).
Their norms are at most \(dmC=d^4m\Delta_K^{1/d}A\).
The case \(m=0\) uses the empty family and bound \(1\).
This proves every clause of the source predicate
\srcid{CompHt.IsShortGen}: equality with the \(\mathbb Z\)-span,
the bound at every embedding, and a bound at least one.
\end{proof}

\subsection{The exact good-cover input packet}
\label{subsec:ftb-good-cover-packet}

For use with the actual theta construction, retain a source core
\(T\) with its laws, an index \(j\) in the prescribed rank interval,
and a good deep cover \(Y\), with coefficient field \(K\) and level
\(M\).  Its numerical bounds are
\[
 p,q\ge1,\quad f>2,\quad L\ge1,\quad
 \log M\le30L,\quad \operatorname{fieldSz}(E)\le3L,\quad
 \log(fr)\le L,\quad o\le391(frL)^2 .
\]
The good-cover bound used below is the single inequality
\[
 \operatorname{fieldSz}(K)
  =\log[K:\mathbb Q]+\frac{\log|\operatorname{disc}K|}{[K:\mathbb Q]}
                                                       \le o .
\]
In particular it bounds each of the two nonnegative summands.
The remaining theta-weight conditions are
\[
 n\ {\rm even},\quad n\ge8fr+4,\quad
 i\longmapsto a_{\tau,i}\text{ is antitone for every }\tau,\quad
 \lambda_{\tau,k}>n/2\ (k<p),\quad
 \log\kappa\le20L .
\]
The source predicate \(\srcid{Dom}(a,n)\) is exactly antitonicity
of every row together with evenness of \(n\); the displayed repetition
of evenness follows the source's input packet.
The good-cover laws retain the same harmonic differential and frame
dictionary as the theta family;
their geometric construction is not a consequence of the inequalities.
This is the full conjunction named \srcid{CovGmin.Hyp} in
\srcid{ThFamMono}.

For clarity, the three substantive statements used by
\srcid{EvenDomainWeightBound.gmin\_good\_cover} have the following
joint quantifier order.  Their constants \(C_A,C_B,C_C\ge0\) precede
the host, rank, core, cover, field, level, row degrees, and \(n\).
For this fixed datum write \(Q=(frL)^2\) and \(\kappa=\operatorname{sizeK}\).
\begin{enumerate}
\item The raw theta family of polynomial norm at most
\(\exp(C_AQ\kappa)\) spans the cusp space over \(K\).
A member comes with a generating level-label tuple \(x_0\), a positive
definite unimodular Hermitian lattice \(\Lambda\), an odd denominator
\(N\), and an integral harmonic polynomial \(P\) of row degrees \(a\).
The lattice is standard at the source's bad integer
\(Q_SqdM\), hence also at \(M\), and satisfies the source level
inclusions at \(N\).
\item For each such lattice there is a nonzero determinant-balancing
scalar \(s\in K\) with
\[
 (s)\prod_\tau\det_\tau(\Lambda\cap\mathcal O_E^n)^{r_\tau}
       \subset \prod_\tau\det_\tau(\Lambda)^{r_\tau},\qquad
 |\sigma(s)|^2 I_\Lambda^r
             \le\exp(\operatorname{fieldSz}(K))^2,\qquad
 I_\Lambda=[\Lambda:\Lambda\cap\mathcal O_E^n].
\]
The raw estimate has the precise premise
\(|\sigma(s)|^2I_\Lambda^r\le B^2\) at every embedding and conclusion
\(\operatorname{homN}(s\,y)\le B\rho\exp(C_BQ\kappa)\), where
\(\rho\) bounds the polynomial norm.
Take \(B=\exp(\operatorname{fieldSz}(K))\).
For that same scalar, every pointwise normalized frame norm of
\(s\,y\) is therefore at most
\[
 \exp(C_AQ\kappa)\exp(C_BQ\kappa)
                          \exp(391Q\kappa).
\]
Here \(\operatorname{fieldSz}(K)\le o\le391Q\) and
\(\kappa\ge1\).  The factor \(I_\Lambda^{r/2}\) in the raw estimate is
cancelled by this same determinant-balancing scalar; a bound on
\(|\sigma(s)|\) alone would not suffice.
\item The single integer \(N_a=Z^2\) of \eqref{eq:ftb-intZ} makes
\(N_a s\,y\) belong to the integral cusp lattice for every raw theta
vector and every determinant-balancing scalar.  Its logarithm is at
most \(C_CQ\kappa\); the explicit denominator proof allows \(C_C=134\).
\end{enumerate}
The first two items are the analytic spanning/norm interface.
The final item is the local-integrality conclusion; its geometric
valuation producers must be proved for the actual model before this
application is unconditional.

\begin{srproposition}[Assembly with the same integer and scalar]
\label{prop:ftb-integral-span-assembly}
Under these three statements there is a subset of the integral cusp
lattice spanning the cusp space over \(K\), all of whose global and
component Petersson norms at every embedding are at most
\[
 \exp\bigl((C_A+C_B+C_C+391)Q\kappa\bigr).
\]
\end{srproposition}
\begin{proof}
The pointwise normalized frame bound implies the same global and
component Petersson bound: the defining integrals are \(L^2\)-norms
against measures of mass at most one.  Points outside the frame
domain contribute zero.  Multiplication by the positive integer
\(N_a\) multiplies each such norm by exactly \(N_a\), at every
embedding.  Since \(N_a\le\exp(C_CQ\kappa)\), the three exponential
factors and the scalar/field factor multiply to the displayed bound.
The resulting vectors are integral by item 3.
For each raw vector \(y\), its replacement is \(N_a s\,y\), a nonzero
scalar multiple.  Replacing every member of a spanning family in this
way preserves its \(K\)-span.  Both the scalar and the integer used in
the norm bound are therefore exactly those used for integrality.
\end{proof}

\paragraph{Source ledger.}
The numerical denominator is
\srcid{IsIntegralFourierCoeffClassIndicatorLevelPoint.intZ}
and \srcid{LogShortThetaSeriesDenominatorBoundLe.intZ\_log\_le}.
The integral assembly is
\srcid{ExistsBoundedIntegerGoodCoverThetaSeriesIntegral.integral\_smul\_thLa}.
The packet and family are \srcid{ThFamMono.Hyp} and \srcid{ThFamMono.thFam}.
The final spanning assembly is
\srcid{EvenDomainWeightBound.gmin\_good\_cover}.
The independent-integer argument is
\srcid{Proofs/CuspFormEvaluations/ExistsShortIndependentIntegers.lean},
using
\srcid{Proofs/CovolumeBounds/FinrankMulDiscrRpowLeSumSqOfMemTraceDual.lean}
and
\srcid{ReciprocityTheory/ReductiveGroups/ArithmeticGeometry/HeightTheory/Transference[A,B].lean}.
The generator construction is
\srcid{ExistsShortGeneratorsOfShortSpanningFamily.exists\_isShortGen\_of\_short\_span}.
The component triangle inequality used when applying it is proved by
writing its good Gram norm as an ordinary Euclidean norm; conjugation
transfers the assertion to the other embedding in each pair
(\srcid{PeterssonNormComponentAddLeOfGramFrame}).

\subsection{The scalar which cancels the lattice index}
\label{subsec:ftb-reflex-small-scalar}

We retain the two determinant lattices
\(\Lambda_n=\Lambda\cap\mathcal O_E^n\subset\Lambda\), and put
\(I=[\Lambda:\Lambda_n]\).  For an \(\mathcal O_E\)-lattice \(L\)
let \(\det L)\) denote its determinant fractional ideal in the fixed
one-dimensional determinant space, and let \(\det_\tau L\) be its
extension to \(\mathcal O_K\) through the coefficient-field embedding
associated with \(\tau:E\hookrightarrow\mathbb C\).

\begin{srlemma}[The exact reflex-ideal scalar]
\label{lem:ftb-reflex-small-scalar}
Suppose \(K\) is the source's coefficient field, hence totally complex,
and the two lattices have finite index \(I>0\).
There is \(s\in K^\times\) such that
\begin{align}
 (s)\prod_\tau(\det_\tau\Lambda_n)^{r_\tau}
       &\subset\prod_\tau(\det_\tau\Lambda)^{r_\tau},
                  \label{eq:ftb-scalar-ideal}\\
 |\sigma(s)|^2 I^r
       &\le\exp(\operatorname{fieldSz}(K))^2
         \quad(\sigma:K\hookrightarrow\mathbb C).
                  \label{eq:ftb-scalar-index}
\end{align}
Here
\(\operatorname{fieldSz}(K)=\log d+\log\Delta_K/d\), where
\(d=[K:\mathbb Q]\) and \(\Delta_K=|\operatorname{disc}K|\).
\end{srlemma}
\begin{proof}
The ideal
\(\mathfrak a=(\det\Lambda_n)(\det\Lambda)^{-1}\) is integral,
and \(N_E\mathfrak a=I\).
To check this directly, localize at every nonzero prime of
\(\mathcal O_E\).  Smith normal form writes the inclusion as
\(\operatorname{diag}(\pi^{e_1},\ldots,\pi^{e_n})\), with \(e_i\ge0\).
Both the determinant-ideal exponent and the length of the quotient are
\(\sum_i e_i\); multiplying residue-field cardinalities gives the norm
identity.  In particular \(\mathfrak a\ne0\).

Let
\[
 \mathfrak B=\prod_{\tau:E\hookrightarrow\mathbb C}
          (\mathfrak a\mathcal O_K)_\tau^{\,r_\tau}.
\]
Then the two type-norm determinant ideals in
\eqref{eq:ftb-scalar-ideal} have quotient \(\mathfrak B\).
Write \(d=2c\).  Extension of ideals has norm
\(N_K((\mathfrak a\mathcal O_K)_\tau)=
 I^{[K:E]}\).  Since
\(\sum_\tau r_\tau=fr\) and \([E:\mathbb Q]=2f\),
\[
                 N_K\mathfrak B=I^{rc}.
\]
These are the norm and determinant calculations in the source's
\srcid{ExistsSmallScalarReflexIdealSelfDualLattice}.

Set \(R=\exp(\operatorname{fieldSz}(K))^2\), and embed \(K\) in
\(\mathbb C^c\) by one embedding from each conjugate pair.
The lattice \(\mathfrak B^{-1}\) has covolume
\(2^{-c}\sqrt{\Delta_K}\,I^{-rc}\).
The product of disks of squared radius \(R/I^r\) has volume
\(\pi^c R^c I^{-rc}\).
It is larger than \(2^{2c}\) times that covolume:
\(\Delta_K\ge1\),
\(\Delta_K\le R^c\), and
\(2^c\sqrt{\Delta_K}<\pi^cR^c\).
Minkowski therefore gives a nonzero
\(s\in\mathfrak B^{-1}\) with
\(|\sigma(s)|^2<R/I^r\) for one embedding from each pair,
and consequently for all embeddings.
This is \eqref{eq:ftb-scalar-index}.
The inclusion \(s\mathfrak B\subset\mathcal O_K\) gives
\eqref{eq:ftb-scalar-ideal}.
Thus the same scalar supplies the determinant-ideal integrality
condition and the index-compensating archimedean inequality.
\end{proof}

\subsection{What nonnegative expansion valuation means}
\label{subsec:ftb-expansion-regularity}

Here is the integral expansion statement used in
\srcid{ReachedSectionRegularOnIntegralModel}; it is separate from the
existence of the polarized Mumford model.
Let \(V\) be a complete DVR of characteristic zero, \(K_0=\operatorname{Frac}V\),
and \(\varpi\) a uniformizer.  Take the actual polarized degeneration
datum and its uniformization over \(V\), with proper abelian quotient
\(B\), a cubical line bundle \(\mathcal L_B\), period lattice \(X\),
period map \(c:X\to B(V)\), and its quadratic valuation \(Q:X\to\mathbb Z\).
For \(k\ge1\) the coefficient bundle indexed by \(\chi\in X\) is
\[
 \mathcal M_{k,\chi}=
  \mathcal L_B^{\,k-1}\otimes t_{c(\chi)}^*\mathcal L_B
       \otimes c(\chi)^*\mathcal L_B^{-1}
       \otimes 0^*\mathcal L_B .
\]
The last two factors are constant line bundles pulled back from \(V\).
This is exactly \(\srcid{fam}\varphi_B(\chi)\) with exponents
\((k-1,1,-1,1)\), not a choice of a new coefficient bundle.

After a valued algebraic extension \(F/K_0\), write
\(V_{\ge b}(t)\) for the following assertion about a coefficient
\(t\in H^0(B_F,\mathcal M_{k,\chi,F})\):
it is a finite sum
\[
                         t=\sum_i a_i t_i
\]
of scalar extensions of integral sections \(t_i\), with
\(v(a_i)\ge b\) for every \(i\).
For a quasi-periodic family \(\Theta=(\Theta_\chi)_{\chi\in X}\),
the source condition \(\nu(\Theta)\ge r\) means
\begin{equation}\label{eq:ftb-nuge-definition}
 V_{\ge r+Q(\chi)/k}(\Theta_\chi)
                          \quad\hbox{for every }\chi\in X.
\end{equation}
The valuation is normalized by \(v(\varpi)=1\).
This definition records finite linear combinations of integral
sections; it is stronger and more precise than a bound on chosen
scalar Fourier coefficients.

The uniformization data used below includes its expansion
correspondence: a generic section admits an expansion after clearing
one nonzero element of \(V\); the integral expansion of such a cleared
section agrees on every thickening modulo \(\varpi^m\) with a locally
finite sum of its coefficient sections.  The coefficient sums and
line-bundle identifications are those of this same uniformization.
This is the part of the actual model construction used here.
The proof below does not manufacture that correspondence from an
arbitrary formal series.

\begin{srlemma}[A finite valued-coefficient retraction]
\label{lem:ftb-valued-retraction}
Let \(F/K_0\) carry an extension of the normalized valuation.
If \(a_1,\ldots,a_t\in F\) have nonnegative valuation, there is a
\(K_0\)-linear map \(\ell:F\to K_0\) with
\(\ell(1)=1\) and \(\ell(a_i)\in V\) for every \(i\).
\end{srlemma}
\begin{proof}
The finite \(V\)-module
\(M=V\cdot1+\sum_iV a_i\) is torsion free, hence free.
Every member of \(M\) has nonnegative valuation.
The coordinate vector of \(1\) in a \(V\)-basis of \(M\) has a unit
coordinate.  Otherwise \(1=\varpi m\) for some \(m\in M\), contradicting
\(0=v(1)=1+v(m)\).
Choose such a coordinate functional and divide it by its value at
\(1\).  It maps \(M\) into \(V\) and sends \(1\) to \(1\).
The chosen \(V\)-basis remains independent over \(K_0\), so extend this
functional from its \(K_0\)-span to \(F\).
This proves the assertion.
\end{proof}

\begin{srlemma}[Descent of a nonnegative coefficient]
\label{lem:ftb-coefficient-descent}
Let \(t\in H^0(B_{K_0},\mathcal M_{k,\chi,K_0})\).
If its scalar extension to \(F\) satisfies \(V_{\ge0}\), then \(t\)
is the generic restriction of an integral section of
\(\mathcal M_{k,\chi}\).
\end{srlemma}
\begin{proof}
Choose a finite representation
\(t_F=\sum_i a_i(t_i)_F\) with \(v(a_i)\ge0\) and integral \(t_i\).
Apply Lemma~\ref{lem:ftb-valued-retraction} to the \(a_i\).
Scalar extension of sections commutes with the coefficient functional:
on a finite affine trivializing cover it applies coefficientwise to
the same gluing equations.  Since \(\ell(1)=1\), this gives
\[
                        t=\sum_i\ell(a_i)t_i .
\]
Every coefficient belongs to \(V\), so the sum is integral.
This is the finite-coefficient argument in
\srcid{SectionRegularOnModelOfNuNonneg.vge\_descend}.
\end{proof}

\begin{srlemma}[Divisibility detected on the same thickening]
\label{lem:ftb-expansion-divisibility}
Suppose an integral section \(\sigma\) of the uniformized line bundle
has its integral expansion on the thickening modulo \(\varpi^m\),
and every coefficient section is divisible by \(\varpi^m\).
Then \(\sigma\) is divisible by \(\varpi^m\) as an integral section.
\end{srlemma}
\begin{proof}
On a local line-bundle trivialization of that thickening, the expansion
formula is a finite sum, each term being a coefficient section times
the prescribed invertible transition factor.  Each coefficient
vanishes modulo \(\varpi^m\), so the restriction of \(\sigma\) is zero.
The underlying group scheme is flat over \(V\).  Thus multiplication
by \(\varpi^m\) on its line bundle is injective, and the kernel of
restriction to the thickening is exactly its image.
The local quotients glue uniquely, yielding the desired section.
This is the argument of
\srcid{JetHarmonics/ExistsDivOfExpandsAt.exists\_div\_of\_expandsAt}.
\end{proof}

\begin{srproposition}[The precise expansion-to-regularity implication]
\label{prop:ftb-expansion-regularity}
With this actual uniformization and its expansion correspondence, a
generic section whose expansion over \(F\) satisfies
\(\nu\ge0\) extends as an integral section of the same model.
\end{srproposition}
\begin{proof}
The positive quadratic valuation in the polarized degeneration has
\(Q(\chi)\ge0\).  Hence \eqref{eq:ftb-nuge-definition} implies
\(V_{\ge0}\) for every coefficient.
One may first take the generic expansion over \(K_0\) and then extend
to \(F\); the expansion identities and uniqueness ensure that this is
the expansion under consideration.
Lemma~\ref{lem:ftb-coefficient-descent} makes every coefficient integral
over \(V\).

Choose a common nonzero \(r\in V\) such that \(r\sigma\) is integral
and its integral expansion is \(r\Theta\); write \(r=u\varpi^m\).
Each coefficient of this integral expansion is divisible by
\(\varpi^m\).  Lemma~\ref{lem:ftb-expansion-divisibility} gives an
integral section \(\tau\) with
\(r\sigma=\varpi^m\tau\).
Thus \(\sigma=u^{-1}\tau\) is integral.
The source implements the clearing, coefficient descent, division,
and unit removal in
\srcid{ExistsIntOfVgeZero.main} and
\srcid{SectionRegularOnModelOfNuNonneg.exists\_extension}.
The remaining cartesian line-bundle comparisons simply transport this
integral section to the chart's \(\srcid{Regular}\) predicate.
They do not assert an extension of a new generic section over an
unrelated model.
\end{proof}

\paragraph{Exact boundary to the model construction.}
Proposition~\ref{prop:ftb-expansion-regularity} uses the actual
uniformization's expansion correspondence and its compatible
identifications with the chart's line bundle.
Existence of those data is the model input, and is not proved merely
by the finite-valued-coefficient lemma.
The additional assertion that every \emph{reached} section has
\(\nu\ge0\) requires the generator, normalized CM pullback, and
symmetric-torsion-translation calculations; it is recorded separately
in the source and is not part of the implication proved here.

\paragraph{Source ledger.}
All paths below have prefix
\srcid{src/Development/AbelianVarietyTheory/Proofs/}.
The small scalar is
\srcid{DivisorProofs/ExistsSmallScalarReflexIdealSelfDualLattice.lean}
(blob \srcid{000df27c7a6e8db0a9255e6d5f7871a0a5152c38}).
The valuation/regularity definitions are
\srcid{Definitions/PlaceAdaptation/AdaptStd.lean} (relative instead to
\srcid{src/Development/AbelianVarietyTheory/}).
The finite functional and coefficient descent are
\srcid{RingsOfIntegers/SectionRegularOnModelOfNuNonneg.lean}
(blob \srcid{7df6dd75a9036946ff91ffd6577e266b7a3f923e}).
Its integral coefficient step is
\srcid{RingsOfIntegers/ExistsIntOfVgeZero.lean}
(blob \srcid{13a90482f19569314e6b1ce248dcd6923f26644e});
the thickening division is
\srcid{JetHarmonics/ExistsDivOfExpandsAt.lean}
(blob \srcid{42e20e8aaac619d94d1f8a27afd63b01ff0551a2}).
The coefficient bundles and expansion relation are defined in
\srcid{src/ReciprocityTheory/ReductiveGroups/ArithmeticGeometry/ExtensionDescent/DegExp.lean};
\(\nu\) is defined in the corresponding
\srcid{DiscreteValuationRings/ValuationMultiplicationInvariance.lean}.
